\section{Inner Monologue: de Open Loop a Feedback-Conditioned Replanning}

SayCan vuelve a elegir una habilidad en cada paso, pero puede no saber si el paso anterior tuvo éxito. Si el robot no logró tomar la lata de refresco de cola y aun así continúa con «llevársela al usuario», el plan de alto nivel es coherente en el lenguaje, pero ya no es válido en el mundo. Inner Monologue añade al planner scene description, clarification y success detector, de modo que el razonamiento en lenguaje cierre realmente el lazo.

\subsection{Agenda y Open-Loop Failure}

La agenda pasa a Inner Monologue. La siguiente slide ilustra el problema con una trayectoria fallida: el usuario pide una bebida y el planner elige el refresco de cola; tras fallar la toma del objeto, el sistema sigue la secuencia original, se desplaza y declara la tarea terminada. Sin observation feedback, el CoT es solo un open-loop script.

\lecturefigure{slide-30-agenda-inner-monologue.jpg}{Agenda: de SayCan a Inner Monologue con retroalimentación del entorno}{Diapositiva V030 recuperada de la grabación oficial; Stanford CS25 Lecture 15}

\lecturefigure{slide-31-saycan-open-loop-failure.jpg}{Open-loop failure: después de fallar la toma del refresco de cola, el planner sigue ejecutando el plan anterior}{Diapositiva V031 recuperada de la grabación oficial; Huang et al., 2022}

\begin{importantbox}{El planning corporeizado debe mantener un estado actualizable}
Un plan en lenguaje lleva implícita una belief: qué hay sobre la mesa, qué tiene el robot en la mano, cuáles son las preferencias del usuario y si el subobjetivo actual ya se cumplió. Cualquier action puede modificar esa belief. Si el sistema no lee las nuevas observaciones, por razonables que sean los tokens siguientes, solo serán un razonamiento sobre un estado del mundo caducado.
\end{importantbox}

\subsection{Passive Scene Description}

La trayectoria fallida de la subsección anterior muestra que un open-loop planner, aunque genere pasos razonables, sigue avanzando sobre un plan caducado después de una sola toma fallida. La corrección más directa no es modificar de inmediato el planner, sino hacer que el estado del entorno regrese de forma continua al contexto en lenguaje. Esta sección examina cómo la passive feedback comprime en texto los resultados de detección y el avance de la tarea; al leer la figura hay que atender tanto a qué estados aporta como a qué hechos físicos se descartan porque la detector ontology es insuficiente.

La retroalimentación pasiva la escribe de forma continua un detector o un VLM, por ejemplo: «veo una lata de refresco de cola y una botella de refresco de lima», «el bloque amarillo ya está en el tazón amarillo». Estas descripciones comprimen imágenes de alta dimensión en un estado textual que el planner puede consumir, de modo que el LLM sabe después de cada paso si el mundo cambió como se esperaba.

\lecturefigure{slide-32-inner-monologue-passive-feedback.jpg}{Passive feedback: object detection, task-progress description y visual QA}{Diapositiva V032 recuperada de la grabación oficial; Huang et al., 2022}

\begin{knowledgebox}{El cuello de botella de información de la passive feedback}
El scene caption conserva solo los conceptos que el detector sabe expresar. Si el sistema no tiene etiquetas como «vaso volcado», «derrame de líquido» o «pinza vacía», el planner no ve esos hechos. Convertir a texto reduce la complejidad de la interfaz, pero también convierte la perception ontology en un límite superior.
\end{knowledgebox}

\subsection{Active Scene Description y Clarification}

La retroalimentación activa se solicita solo cuando el planner tiene incertidumbre, por ejemplo para preguntar por las preferencias del usuario, aclarar una instrucción o consultar la ubicación de un objeto. La slide 33 presenta juntas la active scene description y la visual QA in context: el robot puede preguntar «¿qué bebida quieres?» y actualizar el plan según la respuesta, en lugar de adivinar en condiciones de incertidumbre.

\lecturefigure{slide-33-inner-monologue-active-feedback.jpg}{Active feedback: preguntar al usuario, al entorno o al módulo de preguntas y respuestas visuales solo cuando hace falta}{Diapositiva V033 recuperada de la grabación oficial; Huang et al., 2022}

Si la belief textual se denota $b_t$, la observación de la ejecución $o_{t+1}$, la success signal $z_{t+1}$ y la retroalimentación del usuario $u_{t+1}$, la actualización en lazo cerrado puede abstraerse como
\begin{equation}
b_{t+1}=U(b_t,o_{t+1},z_{t+1},u_{t+1}),
\qquad
k_{t+1}=\arg\max_k \operatorname{Score}(k\mid g,b_{t+1}).
\end{equation}
A diferencia de un open-loop plan, cada elección se basa en la belief actualizada.

\subsection{Success Detector y tabla de resultados}

La mitad superior de la slide de resultados compara combinaciones como CLIPort, oracle, LLM, object feedback e Inner Monologue; la mitad inferior muestra diálogos concretos y el task-progress detector. Al leer la tabla conviene distinguir primero entre seen / unseen tasks y después comparar las diferencias entre añadir solo object feedback, solo success feedback y ambas combinadas.

\lecturefigure{slide-34-inner-monologue-results.jpg}{Inner Monologue results: object feedback, success detector y replanning en lazo cerrado}{Diapositiva V034 recuperada de la grabación oficial; Huang et al., 2022}

\begin{warningbox}{La feedback no es verdadera por sí misma}
El success detector puede dar falsos positivos, la scene description puede omitir objetos y la respuesta del usuario también puede ser ambigua. El closed loop solo da a los errores la oportunidad de corregirse; si el modelo de retroalimentación tiene un sesgo sistemático, el planner seguirá siendo coherente consigo mismo sobre una belief errónea.
\end{warningbox}

\subsection{Los límites de Importing Common Sense}

La última slide contrasta la afirmación de 2018 de que «el cuello de botella de los robots es el high-level semantic planning» con el «no» de los LLM en 2022. La lesson de la parte inferior es más precisa: si el language es la universal API del sistema, los foundation models pueden aportar common sense. Esto reduce el costo de la planificación semántica de alto nivel, pero no elimina los perception, skill y control bottlenecks.

\lecturefigure{slide-35-foundation-common-sense-lesson.jpg}{Aprovechar los foundation models para aportar common sense, sin que sustituyan al sistema físico}{Diapositiva V035 recuperada de la grabación oficial; Stanford CS25 Lecture 15}

\teachervoice{El ponente no presenta al LLM como el «cerebro» del robot. Su argumento es que la language interface permite que el common sense de un modelo externo entre en un sistema existente; si el detector, la affordance o el controller no ofrecen la interfaz correcta, por potente que sea el LLM, no tendrá un estado del mundo sobre el cual operar.}

\begin{lstlisting}[caption={Closed-loop planner de SayCan + Inner Monologue},label={lst:closed-loop-planner}]
belief = describe_scene(initial_observation)
history = []

while not task_complete(belief):
    candidates = language_model.propose_skills(goal, belief, history)
    scored = []
    for skill in candidates:
        useful = language_model.score(skill, goal, history)
        possible = affordance_model.score(skill, belief)
        scored.append((useful * possible, skill))

    skill = max(scored)[1]
    outcome = execute(skill)
    feedback = collect_scene_success_and_user_feedback(outcome)
    belief = update_belief(belief, feedback)
    history.append((skill, feedback))
\end{lstlisting}

\subsection{Resumen de la sección}

Inner Monologue convierte el planning en lenguaje de una open-loop sequence en un observation-conditioned loop. La passive scene, la active clarification y la success detection aportan distintos canales de retroalimentación; la capacidad del sistema sigue limitada por la perception ontology y la detector reliability, pero al menos puede volver a planificar cuando el mundo se aparta del plan.
